% !TEX root = ../main.tex
% ---------------------------------------------------------------------
%  P2 -- SIMD mit ARM NEON
% ---------------------------------------------------------------------
\vlsec{cP2}{P2}{SIMD: ARM NEON}
\label{sec:neon}

\SW{Motivation} Colour adjust, 1600$\times$1200 $\approx$ 2\,Mpixel,
3 Additionen pro Pixel:
\begin{ztab}{@{}Xr@{}}
Variante & Instr./Pixel \\ \midrule
naiv: 3 ADD $+$ LD $+$ ST $+$ DEC $+$ BR & $\geq7$ \\
$+$ Load/Store-Multiple, Unrolling 8 (DEC$+$BR geteilt) & $\geq5.25$ \\
$+$ SIMD 8 & $\geq3.25$ \\
\end{ztab}
Rechnung zu 5.25: pro 8 Pixel $24+8+8+1+1=42$ Instruktionen.
Nach SIMD bleiben im Wesentlichen \emph{Load und Store} übrig
\ra \T{Bottleneck RAM}. 6\,MB Bilddaten passen nicht in 2\,MB L3:
$\sim$95\,000 Misses $\times$ 100 Cycles obendrauf.

\subsection{Lanes}
128-Bit-Register, je nach Befehlssuffix geteilt in
$16\times8$ / $8\times16$ / $4\times32$ / $2\times64$ Bit.
\T{Lane 0 $=$ LSBs}, ganz rechts.
Hardware: Carry-Kette an den Lane-Grenzen auftrennen -- fast gratis,
deshalb ist SIMD so attraktiv.

\begin{formel}[Merksatz]
Der Speedup skaliert mit der Lane-Zahl \ra immer den \emph{kleinsten
ausreichenden} Datentyp wählen. 8\,Bit gibt 16 Lanes, 32\,Bit nur 4.
Darum INT8 in Bildverarbeitung und KI.
\end{formel}

\subsection{Registerfile}
\SW{ARMv7} $16\times$\T{Q} (128\,Bit), $32\times$\T{D} (64),
$32\times$\T{S} (32, VFP, teilt sich Q0--Q7 / D0--D15).
Alle auf denselben Flip-Flops -- \emph{wie eine C-Union}.
\SW{AArch64} (Pi 4/5 mit 64-Bit-OS) $32\times$\T{V}, 128\,Bit, kein
Paar-Aliasing, Syntax \R{ADD V0.4S,V1.4S,V2.4S}.

\begin{falle}[Zwei Fallen beim Registerfile]
1.~Schreiben in \R{D3} zerstört die obere Hälfte von \R{Q1}.\\
2.~NEON hat ein \emph{eigenes} Registerfile und eine eigene Pipeline
\ra Transfers ARM\,$\leftrightarrow$\,NEON kosten viele Cycles. Code, der
ständig hin- und herspringt, ist \emph{langsamer} als reiner Skalarcode.
\end{falle}

\subsection{Datentypen}
\begin{ztab}{@{}lXX@{}}
 & 8/16/32/64 & Bemerkung \\ \midrule
U / S & \R{U8}\dots\R{S64} & unsigned / signed \\
I & \R{I8}\dots\R{I64} & vorzeichen\emph{un}abhängig \\
F & \R{F16}, \R{F32} & \T{kein F64} in NEON! \\
P & \R{P8}, \R{P16} & Polynom über GF(2), carry-less \\
\end{ztab}
\SW{Warum „I“} Addition ist bitweise identisch für signed/unsigned
\ra \R{VADD.I16}. Sättigung ist es \emph{nicht}
\ra nur \R{VQADD.S16}/\R{.U16}. Die Existenz von „I“ verrät also, dass der
Befehl vorzeichenunabhängig ist.
\SW{P} für CRC, AES-GCM, Reed-Solomon (XOR statt Addition, kein Übertrag).
\SW{Float} NEON-Floats sind auf ARMv7 nicht voll IEEE-754 (flush-to-zero)
\ra GCC vektorisiert Float-Schleifen nur mit \R{-ffast-math}.

\subsection{Befehlssyntax dekodieren}
\begin{formel}[Aufbau jedes NEON-Befehls]
\label{rez:neondecode}
\begin{center}\ttfamily\small
V\{mod\}op\{shape\}\{cond\}.\{dt\}\ dest,\ src1,\ src2
\end{center}
\begin{tabularx}{\linewidth}{@{}clX@{}}
\multicolumn{3}{@{}l}{\T{mod}}\\
Q & saturating & klemmt statt umzuwickeln ($250{+}10=255$) \\
H & halving & $(a{+}b)/2$ ohne Überlauf \\
D & doubling & Q15-Festkomma: 1 Bit nach links \\
R & rounding & statt abschneiden, vermeidet Bias \\[2pt]
\multicolumn{3}{@{}l}{\T{shape} (Ziel $\neq$ Quellformat)}\\
-- & & D,\,D \ra D \\
L & long & D,\,D \ra Q \\
W & wide & Q,\,D \ra Q \\
N & narrow & Q,\,Q \ra D \\
\end{tabularx}
\end{formel}

\begin{bsp}[Beispiel: \R{VQRDMULH.S16 Q0,Q1,Q2}]
V$|$\,Q\,R\,D\,$|$MULH$|$.S16\\
$=$ Saturating $+$ Rounding $+$ Doubling Multiply, High half,
8 Lanes à 16\,Bit signed. Der Arbeitsbefehl jedes FIR-Filters --
gelesen, ohne ihn zu kennen.
\end{bsp}

\SW{Standardmuster} 8\,Bit laden \ra mit \T{L} auf 16\,Bit erweitern
\ra in 16\,Bit rechnen \ra mit \T{N} und Sättigung zurück auf 8\,Bit.

\subsection{Befehlsklassen}
\begin{description}
\item[Load/Store] \R{VLDn}/\R{VSTn} ($n{=}1..4$) laden
  \emph{de-interleaved}: \R{VLD3.8} liest 24\,B und sortiert R/G/B in drei
  Register. \emph{Ohne das wäre NEON für Bildverarbeitung nutzlos.}
  Dazu \R{VLDM/VSTM}, \R{VLDR/VSTR}, \R{VPUSH/VPOP}.
\item[Data Processing] \R{VDUP} (Skalar in alle Lanes),
  \R{VMOV} (Immediate), \R{VZIP/VUZP}, \R{VTRN} ($4{\times}4$
  transponieren), \R{VTBL/VTBX} (beliebige Permutation),
  \R{VEXT} (gleitendes Fenster über zwei Register),
  \R{VCVT} (auch Fixed-Point), \R{VREV16}.
\item[Shift] Es gibt \emph{keine} SIMD-Division, darum sind Shifts
  zentral: \R{VSHL/VSHR}, \R{VSRA} (shift $+$ accumulate),
  \R{VSHLL}, \R{VQRSHRN} (Sat$+$Round$+$Shift$+$Narrow -- der
  Standard-Abschlussbefehl), \R{VSLI/VSRI} (Bitfelder packen).
\item[Logical/Compare] Prädikation statt Branch:
  \R{VCEQ/VCGE/VCGT} liefern eine \emph{Maske} (alle 1 / alle 0),
  \R{VBSL} (Bit Select) wählt bitweise \ra das \R{?:} der Vektorwelt.
  Preis: beide Zweige werden gerechnet.
\item[Arithmetik] \R{VABD/VABA} (SAD -- Bewegungsschätzung im Videocodec),
  \R{VMAX/VMIN}, \R{VPADD/VPMAX} (\emph{pairwise} $=$ benachbarte Lanes),
  \R{VCLZ/VCNT}, \R{VRECPE/VRECPS} und \R{VRSQRTE/VRSQRTS}
  (Newton-Raphson statt Division/Wurzel).
\item[Multiply] \R{VMLA/VMLS} (MAC -- Grundoperation von FIR,
  Skalarprodukt, Faltung, Matrixmultiplikation), \R{VFMA}
  (\emph{fused}: nur eine Rundung, numerisch genauer),
  „by scalar“-Varianten, \R{VQDMULH} (Q15).
\end{description}

\subsection{Intrinsics}
\begin{ztab}{@{}lX@{}}
\R{v<op>q\_<dt>} & \R{q} $=$ Quadword (128\,Bit), ohne \R{q} $=$ 64 \\
\R{<typ><br>x<anz>\_t} & \R{uint8x16\_t}, \R{float32x4\_t} \\
\R{...x<länge>\_t} & \R{uint8x8x3\_t} $=$ 3 Register (für \R{VLD3}) \\
\end{ztab}
Header \R{<arm\_neon.h>}. Registerzuteilung und Scheduling macht weiterhin
der Compiler.

\SW{Matrixmultiplikation $4\times4$} Daten \T{column-major}
(\R{C[4*j+i]}). Nicht Skalarprodukt Zeile$\times$Spalte (bräuchte
horizontale Summen), sondern jede Spalte von C als \T{Linearkombination
der Spalten von A}: \R{C0 += A0*B0[0] + A1*B0[1] + \dots} über
\R{vfmaq\_laneq\_f32}. \ra Alle 4 Lanes arbeiten durchgehend, \emph{keine}
horizontale Operation, keine Transposition. 12 von 16 Q-Registern belegt --
Registerdruck ist die harte Grenze ($8\times8$ ginge nicht mehr).

\SW{MAC-/FIR-Filter} In der Schleife nur \emph{vertikal} rechnen
(\R{vmla.f32 q0,q1,q2} lane-weise; \R{!} $=$ Post-Increment spart die
Adressrechnung). Am \emph{Ende} einmal horizontal reduzieren:
\R{vpadd} in $\log_2 n$ Schritten
$[a,b],[c,d]\to[a{+}b,\,c{+}d]\to$ Gesamtsumme.

\begin{falle}[Horizontal nie in der Schleife]
Pairwise/horizontale Operationen brechen die Lane-Isolation und sind
langsam. Ein Akkumulator \emph{pro Lane}, Reduktion genau einmal am Schluss.
\end{falle}

\subsection{Grenzen}
\label{sec:simdgrenzen}
\begin{formel}[Amdahl für SIMD]
\label{form:amdahl}
\[S=\frac{1}{(1-p)+p/s}\qquad p=\text{vektorisierbar},\ s=\text{Lanes}\]
$p{=}0.8,\ s{=}4\Rightarrow2.5\times$; $s{=}16\Rightarrow$ nur $4\times$.
Die 20\,\% Skalarcode deckeln alles.
\end{formel}

\begin{description}
\item[Memory bound] \R{c[i]=a[i]+b[i]}: 1 Rechenoperation auf 3
  Speicherzugriffe. \emph{Zuerst} Datenlayout und Cache-Lokalität,
  \emph{dann} SIMD.
\item[Hindernisse] loop-carried dependency (\R{a[i]=a[i-1]+b[i]});
  Pointer-Aliasing \ra \R{restrict}; Kontrollfluss \ra Masken;
  Float-Assoziativität \ra \R{-ffast-math}.
\item[Tail] 1003 Elemente bei 4 Lanes $=$ 250 Vektoriterationen $+$ 3
  skalar. SIMD lohnt erst ab einigen Dutzend Elementen.
\item[Alignment] unit stride ist Pflicht; Gather zerstört die Effizienz,
  weil jede Lane eine eigene Cache Line anfasst.
\end{description}

\subsection{NEON einschalten (bare metal)}
Nach Reset aus -- wegen Strom (clock gating) und lazy Context Switch des OS.
\T{Zwei} Schalter nötig:
\begin{lstlisting}
MRC p15,0,r0,c1,c0,2   ; CPACR lesen
ORR r0,r0,#0x00f00000  ; cp10+cp11 full access
MCR p15,0,r0,c1,c0,2   ; CPACR schreiben
ISB                    ; Pipeline neu holen!
MOV r0,#0x40000000     ; Bit30 = EN
VMSR FPEXC,r0          ; Einheit selbst einschalten
\end{lstlisting}
CPACR ist die Tür des ARM-Kerns zum Coprozessor-Interface,
\R{FPEXC.EN} der Hauptschalter der Einheit. Ohne \R{ISB} ist das
Verhalten undefiniert.
\SW{Flags} \R{armcc -O3 -Otime --vectorize} (veraltet) bzw.
\R{gcc -O3 -ftree-vectorize}; Diagnose \R{-fopt-info-vec-missed}.
